% answers.tex —— 2014 年·陕西卷·理科 · 参考答案与详细解析片段
% 由 tools/generate_projects.py 依来源总账自 2014/2014_数学真题_陕西理_试题.md + 2014/2014_数学真题_陕西理_解析正文.md + 2014/2014_数学真题_陕西理_答案要点.md 生成，勿手改（改源 MD 后重跑）


\begin{answer}
\gksolhead{1}{}

\textbf{题目}：

已知集合 $M=\{x\mid x\geqslant 0,x\in\mathbf R\}$，$N=\{x\mid x^2<1,x\in\mathbf R\}$，则 $M\cap N=$\mbox{（　　）}

（A）$[0,1]$；（B）$[0,1)$；（C）$(0,1]$；（D）$(0,1)$。

\Needspace{4\baselineskip}
\textbf{解析}

\Needspace{5\baselineskip}
先化简两个集合。由 $x^2<1$ 得
\nopagebreak[4]
\[-1<x<1,\]
\Needspace{5\baselineskip}
故
\nopagebreak[4]
\[N=(-1,1).\]
\Needspace{5\baselineskip}
又
\nopagebreak[4]
\[M=[0,+\infty).\]
\Needspace{5\baselineskip}
在数轴上取公共部分：左端点 $0$ 同时属于 $M$ 和 $N$，右端点 $1$ 不属于 $N$，所以
\nopagebreak[4]
\[M\cap N=[0,1).\]

选 B。

\gkanswer{答案：B。$N=(-1,1)$，与 $M=[0,+\infty)$ 取交集得 $[0,1)$。}

\end{answer}
\begin{answer}
\gksolhead{2}{}

\textbf{题目}：

函数 $f(x)=\cos\left(2x-\dfrac\pi6\right)$ 的最小正周期是\mbox{（　　）}

（A）$\dfrac\pi2$；（B）$\pi$；（C）$2\pi$；（D）$4\pi$。

\Needspace{4\baselineskip}
\textbf{解析}

\Needspace{5\baselineskip}
函数 $y=\cos(\omega x+\varphi)$ 的最小正周期为
\nopagebreak[4]
\[T=\frac{2\pi}{|\omega|}.\]
\Needspace{5\baselineskip}
这里 $\omega=2$，故
\nopagebreak[4]
\[T=\frac{2\pi}{2}=\pi.\]

选 B。

\gkanswer{答案：B。$\cos(2x-\dfrac\pi6)$ 的周期为 $\dfrac{2\pi}{2}=\pi$。}

\end{answer}
\begin{answer}
\gksolhead{3}{}

\textbf{题目}：

定积分 $\displaystyle\int_0^1(2x+\mathrm e^x)\,\mathrm dx$ 的值为\mbox{（　　）}

（A）$\mathrm e+2$；（B）$\mathrm e+1$；（C）$\mathrm e$；（D）$\mathrm e-1$。

\Needspace{4\baselineskip}
\textbf{解析}

\Needspace{5\baselineskip}
被积函数的一个原函数为 $x^2+\mathrm e^x$，由微积分基本定理
\nopagebreak[4]
\[\begin{aligned}
\int_0^1(2x+\mathrm e^x)\,\mathrm dx
&=\left[x^2+\mathrm e^x\right]_0^1\\
&=(1+\mathrm e)-(0+1)=\mathrm e.
\end{aligned}\]

选 C。

\gkanswer{答案：C。$\displaystyle\int_0^1(2x+\mathrm e^x)\,\mathrm dx=\left[x^2+\mathrm e^x\right]_0^1=\mathrm e$。}

\end{answer}
\begin{answer}
\gksolhead{4}{}

\textbf{题目}：

根据如图所示的框图，对大于 $2$ 的整数 $N$，输出的数列的通项公式是\mbox{（　　）}
\nopagebreak[4]
\begin{center}% ============================================================
% f2014_shaanxi_science_q04_flowchart.tex
% 2014 陕西卷（理科）第 4 题：程序框图
% 来源 PDF：普通高考/2014/2014陕西理.pdf
% 原页：第 1 页（A3），印面 911，第 4 题题干下方
% 按原图重绘；判断框内为“$i>N$”，否分支回到“$a_i=2*S$”。
% ============================================================
\begin{tikzpicture}[>=Stealth, line width=0.7pt, font=\small,
  every node/.style={inner sep=2pt},
  term/.style={draw, rounded corners=2pt, minimum width=2.35cm,
    minimum height=0.62cm, align=center},
  proc/.style={draw, minimum width=2.75cm, minimum height=0.62cm, align=center},
  dec/.style={draw, diamond, aspect=1.9, minimum width=2.35cm,
    minimum height=0.95cm, align=center}]
  \node[term] (start) at (0,0) {开始};
  \node[proc] (input) at (0,-0.95) {输入 $N$};
  \node[proc] (init) at (0,-1.9) {$S=1,\ i=1$};
  \node[proc] (body) at (0,-2.85) {$a_i=2*S$};
  \node[proc] (setS) at (0,-3.8) {$S=a_i$};
  \node[proc] (inc) at (0,-4.75) {$i=i+1$};
  \node[dec] (test) at (0,-5.95) {$i>N$};
  \node[proc, minimum width=3.55cm] (out) at (0,-7.1) {输出 $a_1,a_2,\cdots,a_N$};
  \node[term] (end) at (0,-8.05) {结束};

  \draw[->] (start) -- (input);
  \draw[->] (input) -- (init);
  \draw[->] (init) -- (body);
  \draw[->] (body) -- (setS);
  \draw[->] (setS) -- (inc);
  \draw[->] (inc) -- (test);
  \draw[->] (test) -- node[right] {是} (out);
  \draw[->] (out) -- (end);
  % “否”分支：从判断框左侧绕回循环体
  \draw[->] (test.west) -- ++(-1.55,0) |- (body.west);
  \node at (-1.05,-5.72) {否};
\end{tikzpicture}
\end{center}
\nopagebreak[4]
（A）$a_n=2n$；（B）$a_n=2(n-1)$；（C）$a_n=2^n$；（D）$a_n=2^{n-1}$。

\Needspace{4\baselineskip}
\textbf{解析}

按框图逐步执行。初始 $S=1$，$i=1$。

\Needspace{5\baselineskip}
第一次循环：
\nopagebreak[4]
\[a_1=2S=2,\qquad S=a_1=2,\qquad i=2.\]
\Needspace{5\baselineskip}
第二次循环：
\nopagebreak[4]
\[a_2=2S=4,\qquad S=a_2=4,\qquad i=3.\]
\Needspace{5\baselineskip}
第三次循环：
\nopagebreak[4]
\[a_3=2S=8,\qquad S=a_3=8,\qquad i=4.\]
\Needspace{5\baselineskip}
可见每循环一次，$a_i$ 在上一项基础上乘以 $2$，且 $a_1=2$，故
\nopagebreak[4]
\[a_n=2\cdot2^{n-1}=2^n.\]
当 $i>N$ 时输出前 $N$ 项，所以通项公式为 $a_n=2^n$。

选 C。

\gkanswer{答案：C。循环依次得到 $a_1=2$，$a_2=4$，$a_3=8$，故 $a_n=2^n$。}

\end{answer}
\begin{answer}
\gksolhead{5}{}

\textbf{题目}：

已知底面边长为 $1$，侧棱长为 $\sqrt2$ 的正四棱柱的各顶点均在同一个球面上，则该球的体积为\mbox{（　　）}

（A）$\dfrac{32\pi}{3}$；（B）$4\pi$；（C）$2\pi$；（D）$\dfrac{4\pi}{3}$。

\Needspace{4\baselineskip}
\textbf{解析}

正四棱柱的底面是边长为 $1$ 的正方形，侧棱长为 $\sqrt2$。它的八个顶点都在同一个球面上，该球以正四棱柱的体对角线为直径。

\Needspace{5\baselineskip}
体对角线长为
\nopagebreak[4]
\[d=\sqrt{1^2+1^2+(\sqrt2)^2}=\sqrt{1+1+2}=2.\]
\Needspace{5\baselineskip}
故球半径
\nopagebreak[4]
\[R=\frac d2=1.\]
\Needspace{5\baselineskip}
球的体积为
\nopagebreak[4]
\[V=\frac43\pi R^3=\frac43\pi.\]

选 D。

\gkanswer{答案：D。正四棱柱的体对角线为球的直径，$2R=\sqrt{1^2+1^2+(\sqrt2)^2}=2$，故 $R=1$，体积为 $\dfrac{4\pi}{3}$。}

\end{answer}
\begin{answer}
\gksolhead{6}{}

\textbf{题目}：

从正方形四个顶点及其中心这 $5$ 个点中，任取 $2$ 个点，则这 $2$ 个点的距离不小于该正方形边长的概率为\mbox{（　　）}

（A）$\dfrac15$；（B）$\dfrac25$；（C）$\dfrac35$；（D）$\dfrac45$。

\Needspace{4\baselineskip}
\textbf{解析}

\Needspace{5\baselineskip}
$5$ 个点为正方形的四个顶点 $A,B,C,D$ 和中心 $O$。任取 $2$ 点，共有
\nopagebreak[4]
\[\mathrm C_5^2=10\]
种取法。

按距离分类：

\begin{itemize}
\item 四个顶点两两配对共 $\mathrm C_4^2=6$ 对，距离为边长或对角线长，均不小于边长；
\item 中心 $O$ 与四个顶点配对共 $4$ 对，距离为 $\dfrac{\sqrt2}{2}$ 倍边长，小于边长。
\end{itemize}

\Needspace{5\baselineskip}
所以符合条件的取法有 $6$ 种，所求概率为
\nopagebreak[4]
\[P=\frac6{10}=\frac35.\]

选 C。

\gkanswer{答案：C。$5$ 个点任取 $2$ 点共 $\mathrm C_5^2=10$ 种，其中距离不小于边长的有 $4$ 条边和 $2$ 条对角线，共 $6$ 种，概率为 $\dfrac35$。}

\end{answer}
\begin{answer}
\gksolhead{7}{}

\textbf{题目}：

下列函数中，满足“$f(x+y)=f(x)f(y)$”的单调递增函数是\mbox{（　　）}

（A）$f(x)=x^{\frac12}$；（B）$f(x)=x^3$；（C）$f(x)=\left(\dfrac12\right)^x$；（D）$f(x)=3^x$。

\Needspace{4\baselineskip}
\textbf{解析}

条件 $f(x+y)=f(x)f(y)$ 是指数函数的特征性质。逐一检验：

\Needspace{5\baselineskip}
（A）$f(x)=x^{1/2}$：
\nopagebreak[4]
\[f(x+y)=\sqrt{x+y},\qquad f(x)f(y)=\sqrt{xy},\]
一般不相等。

\Needspace{5\baselineskip}
（B）$f(x)=x^3$：
\nopagebreak[4]
\[(x+y)^3\neq x^3y^3,\]
一般不相等。

（C）$f(x)=\left(\dfrac12\right)^x$：满足 $f(x+y)=f(x)f(y)$，但底数 $\dfrac12<1$，函数单调递减。

\Needspace{5\baselineskip}
（D）$f(x)=3^x$：
\nopagebreak[4]
\[f(x+y)=3^{x+y}=3^x\cdot3^y=f(x)f(y),\]
且底数 $3>1$，函数单调递增。

选 D。

\gkanswer{答案：D。$3^{x+y}=3^x\cdot3^y$，且 $y=3^x$ 单调递增。}

\end{answer}
\begin{answer}
\gksolhead{8}{}

\textbf{题目}：

原命题为“若 $z_1$，$z_2$ 互为共轭复数，则 $|z_1|=|z_2|$”，关于其逆命题、否命题、逆否命题真假性的判断依次如下，正确的是\mbox{（　　）}

（A）真，假，真；（B）假，假，真；（C）真，真，假；（D）假，假，假。

\Needspace{4\baselineskip}
\textbf{解析}

\Needspace{5\baselineskip}
设
\nopagebreak[4]
\[p:\ z_1,z_2\text{ \text{互为共轭复数}},\qquad q:\ |z_1|=|z_2|.\]
\Needspace{5\baselineskip}
原命题为 $p\Rightarrow q$。若 $z_2=\bar z_1$，则
\nopagebreak[4]
\[|z_2|=|\bar z_1|=|z_1|,\]
故原命题为真。

\Needspace{5\baselineskip}
逆命题为 $q\Rightarrow p$。取
\nopagebreak[4]
\[z_1=1,\qquad z_2=-1,\]
则 $|z_1|=|z_2|=1$，但 $z_2\neq\bar z_1$，故逆命题为假。

否命题为 $\neg p\Rightarrow\neg q$。仍取上述 $z_1=1$、$z_2=-1$：它们不互为共轭，却有 $|z_1|=|z_2|$，故否命题为假。

逆否命题为 $\neg q\Rightarrow\neg p$，与原命题等价，故为真。

因此逆命题、否命题、逆否命题的真假依次为“假、假、真”。

选 B。

\gkanswer{答案：B。原命题真；逆命题“模相等则互为共轭”假；否命题与原命题同真假，故假；逆否命题与原命题同真假，故真。}

\end{answer}
\begin{answer}
\gksolhead{9}{}

\textbf{题目}：

设样本数据 $x_1,x_2,\cdots,x_{10}$ 的均值和方差分别为 $1$ 和 $4$，若 $y_i=x_i+a$（$a$ 为非零常数，$i=1,2,\cdots,10$），则 $y_1,y_2,\cdots,y_{10}$ 的均值和方差分别为\mbox{（　　）}

（A）$1+a$，$4$；（B）$1+a$，$4+a$；（C）$1$，$4$；（D）$1$，$4+a$。

\Needspace{4\baselineskip}
\textbf{解析}

\Needspace{5\baselineskip}
设原数据的平均数为 $\bar x=1$。新数据 $y_i=x_i+a$ 的平均数为
\nopagebreak[4]
\[\bar y=\frac1{10}\sum_{i=1}^{10}(x_i+a)=\bar x+a=1+a.\]
\Needspace{5\baselineskip}
方差衡量数据相对均值的波动，整体平移不改变波动：
\nopagebreak[4]
\[\begin{aligned}
s_y^2&=\frac1{10}\sum_{i=1}^{10}(y_i-\bar y)^2\\
&=\frac1{10}\sum_{i=1}^{10}(x_i+a-1-a)^2\\
&=\frac1{10}\sum_{i=1}^{10}(x_i-1)^2=s_x^2=4.
\end{aligned}\]
故新数据的均值和方差分别为 $1+a$ 和 $4$。

选 A。

\gkanswer{答案：A。均值随平移变为 $1+a$，方差不变，仍为 $4$。}

\end{answer}
\begin{answer}
\gksolhead{10}{}

\textbf{题目}：

如图，某飞行器在 $4$ 千米高空水平飞行，从距着陆点 $A$ 的水平距离 $10$ 千米处下降，已知下降飞行轨迹为某三次函数图象的一部分，则函数的解析式为\mbox{（　　）}
\nopagebreak[4]
\begin{center}% ============================================================
% f2014_shaanxi_science_q10_cubic.tex
% 2014 陕西卷（理科）第 10 题：三次函数下降轨迹
% 来源 PDF：普通高考/2014/2014陕西理.pdf
% 原页：第 1 页（A3），印面 911，第 10 题题干下方
% 曲线 y=x^3/125-3x/5 在 x=-5 处高 2、在 x=5 处低 -2，
% 两端切线水平；虚线为高度与水平距离的辅助线。
% ============================================================
\begin{tikzpicture}[>=Stealth, line width=0.7pt, font=\small, x=0.90cm, y=1.00cm]
  % 坐标轴
  \draw[->] (-6.4,0) -- (6.6,0) node[right] {$x$};
  \draw[->] (0,-3.1) -- (0,3.15) node[above] {$y$};
  \node[below left] at (0,0) {$O$};
  \node[below left, fill=white, inner sep=1pt] at (-5,0) {$-5$};
  \node[below right, fill=white, inner sep=1pt] at (5,0) {$5$};
  \node[above left, fill=white, inner sep=1pt] at (0,2) {$2$};
  \node[below left, fill=white, inner sep=1pt] at (0,-2) {$-2$};

  % 辅助虚线
  \draw[dashed] (-5,2) -- (0,2);
  \draw[dashed] (-5,2) -- (-5,0);
  \draw[dashed] (5,-2) -- (0,-2);
  \draw[dashed] (5,-2) -- (5,0);

  % 下降轨迹
  \draw[thick, dashed, domain=-5:5, samples=120]
    plot (\x,{(\x)^3/125-3*(\x)/5});

  % 起点箭头
  \draw[->, thick] (-5.55,2.02) -- (-5.12,2.02);

  % 着陆点与地面跑道
  \fill (5,-2) circle (1.2pt);
  \node[below right] at (5,-2) {$A$};
  \draw (-5.9,-2) -- (6.25,-2);
  \node[below] at (3.45,-2.10) {地面跑道};
\end{tikzpicture}
\end{center}
\nopagebreak[4]
（A）$y=\dfrac{1}{125}x^3-\dfrac35x$；（B）$y=\dfrac{2}{125}x^3-\dfrac45x$；（C）$y=\dfrac{3}{125}x^3-x$；（D）$y=-\dfrac{3}{125}x^3+\dfrac15x$。

\Needspace{4\baselineskip}
\textbf{解析}

\Needspace{5\baselineskip}
设三次函数为
\nopagebreak[4]
\[f(x)=Ax^3+Bx^2+Cx+D.\]
\Needspace{5\baselineskip}
由图象过原点 $O(0,0)$，得 $D=0$。图象左端点为 $(-5,2)$，右端点为 $(5,-2)$，且在这两点处轨迹切线水平，即
\nopagebreak[4]
\[f(-5)=2,\qquad f(5)=-2,\qquad f'(-5)=0,\qquad f'(5)=0.\]
\Needspace{5\baselineskip}
求导：
\nopagebreak[4]
\[f'(x)=3Ax^2+2Bx+C.\]
\Needspace{5\baselineskip}
由 $f'(-5)=f'(5)=0$ 得
\nopagebreak[4]
\[75A-10B+C=0,\qquad 75A+10B+C=0,\]
\Needspace{5\baselineskip}
两式相减得 $B=0$，再代回得
\nopagebreak[4]
\[C=-75A.\]
\Needspace{5\baselineskip}
于是
\nopagebreak[4]
\[f(x)=Ax^3-75Ax.\]
\Needspace{5\baselineskip}
由 $f(5)=-2$：
\nopagebreak[4]
\[125A-375A=-250A=-2,\]
\Needspace{5\baselineskip}
故
\nopagebreak[4]
\[A=\frac1{125},\qquad C=-75\cdot\frac1{125}=-\frac35.\]
\Needspace{5\baselineskip}
所以
\nopagebreak[4]
\[f(x)=\frac1{125}x^3-\frac35x.\]
检验 $f(-5)=-\dfrac{125}{125}+\dfrac{15}{5}=-1+3=2$，符合左端点。

选 A。

\gkanswer{答案：A。曲线过 $(-5,2)$、$(0,0)$、$(5,-2)$，且在 $x=\pm5$ 处切线水平，只有 $y=\dfrac{1}{125}x^3-\dfrac35x$ 同时满足。}

\end{answer}
\begin{answer}
\gksolhead{11}{}

\textbf{题目}：

已知 $4^a=2$，$\lg x=a$，则 $x=$\gkblank{}。

\Needspace{4\baselineskip}
\textbf{解析}

\Needspace{5\baselineskip}
由 $4^a=2$，即
\nopagebreak[4]
\[2^{2a}=2^1,\]
\Needspace{5\baselineskip}
得
\nopagebreak[4]
\[2a=1,\qquad a=\frac12.\]
\Needspace{5\baselineskip}
又 $\lg x=a$，即
\nopagebreak[4]
\[\lg x=\frac12,\]
\Needspace{5\baselineskip}
故
\nopagebreak[4]
\[x=10^{1/2}=\sqrt{10}.\]

\gkanswer{答案：$\sqrt{10}$。$4^a=2$ 得 $a=\dfrac12$，故 $x=10^{1/2}=\sqrt{10}$。}

\end{answer}
\begin{answer}
\gksolhead{12}{}

\textbf{题目}：

若圆 $C$ 的半径为 $1$，其圆心与点 $(1,0)$ 关于直线 $y=x$ 对称，则圆 $C$ 的标准方程为\gkblank{}。

\Needspace{4\baselineskip}
\textbf{解析}

\Needspace{5\baselineskip}
点 $(1,0)$ 关于直线 $y=x$ 的对称点，是把横、纵坐标互换，得
\nopagebreak[4]
\[(0,1).\]
\Needspace{5\baselineskip}
这就是圆心。圆半径为 $1$，故圆的标准方程为
\nopagebreak[4]
\[(x-0)^2+(y-1)^2=1,\]
\Needspace{5\baselineskip}
即
\nopagebreak[4]
\[x^2+(y-1)^2=1.\]

\gkanswer{答案：$x^2+(y-1)^2=1$。圆心 $(1,0)$ 关于 $y=x$ 的对称点为 $(0,1)$。}

\end{answer}
\begin{answer}
\gksolhead{13}{}

\textbf{题目}：

设 $0<\theta<\dfrac\pi2$，向量 $\boldsymbol a=(\sin2\theta,\cos\theta)$，$\boldsymbol b=(\cos\theta,1)$，若 $\boldsymbol a\parallel\boldsymbol b$，则 $\tan\theta=$\gkblank{}。

\Needspace{4\baselineskip}
\textbf{解析}

\Needspace{5\baselineskip}
两个向量平行的充要条件是行列式为零：
\nopagebreak[4]
\[\begin{vmatrix}
\sin2\theta & \cos\theta\\
\cos\theta & 1
\end{vmatrix}=0.\]
\Needspace{5\baselineskip}
即
\nopagebreak[4]
\[\sin2\theta-\cos^2\theta=0.\]
\Needspace{5\baselineskip}
用二倍角公式：
\nopagebreak[4]
\[2\sin\theta\cos\theta=\cos^2\theta.\]
\Needspace{5\baselineskip}
因为 $0<\theta<\dfrac\pi2$，所以 $\cos\theta>0$，两边除以 $\cos^2\theta$ 得
\nopagebreak[4]
\[2\tan\theta=1,\]
\Needspace{5\baselineskip}
故
\nopagebreak[4]
\[\tan\theta=\frac12.\]

\gkanswer{答案：$\dfrac12$。由 $\boldsymbol a\parallel\boldsymbol b$ 得 $\sin2\theta=\cos^2\theta$，即 $2\sin\theta\cos\theta=\cos^2\theta$，故 $\tan\theta=\dfrac12$。}

\end{answer}
\begin{answer}
\gksolhead{14}{}

\textbf{题目}：

观察分析下表中的数据：

\begin{center}
\begin{tabular}{llll}
\toprule
多面体 & 面数（$F$） & 顶点数（$V$） & 棱数（$E$） \\
\midrule
三棱锥 & $5$ & $6$ & $9$ \\
五棱锥 & $6$ & $6$ & $10$ \\
立方体 & $6$ & $8$ & $12$ \\
\bottomrule
\end{tabular}
\end{center}

猜想一般凸多面体中，$F$，$V$，$E$ 所满足的等式是\gkblank{}。

\Needspace{4\baselineskip}
\textbf{解析}

把表中三行分别代入 $F+V-E$：

\Needspace{5\baselineskip}
三棱锥一行：
\nopagebreak[4]
\[5+6-9=2;\]
\Needspace{5\baselineskip}
五棱锥一行：
\nopagebreak[4]
\[6+6-10=2;\]
\Needspace{5\baselineskip}
立方体一行：
\nopagebreak[4]
\[6+8-12=2.\]
\Needspace{5\baselineskip}
三行结果相同，因此猜想一般凸多面体满足
\nopagebreak[4]
\[F+V-E=2.\]
这正是欧拉多面体公式。原页首行把面数 $5$、顶点数 $6$、棱数 $9$ 的几何体标作“三棱锥”，该组数据实际对应三棱柱；这里照录原页，不影响 $F+V-E=2$ 的结论。

\gkanswer{答案：$F+V-E=2$。表中三行均满足 $5+6-9=6+6-10=6+8-12=2$。}

\end{answer}
\begin{answer}
\gksolhead{15}{}

\textbf{题目}：

三选一。

【A】设 $a,b,m,n\in\mathbf R$，且 $a^2+b^2=5$，$ma+nb=5$，则 $\sqrt{m^2+n^2}$ 的最小值为\gkblank{}。

【B】如图，$\triangle ABC$ 中，$BC=6$，以 $BC$ 为直径的半圆分别交 $AB$，$AC$ 于点 $E$，$F$，若 $AC=2AE$，则 $EF=$\gkblank{}。
\nopagebreak[4]
\begin{center}% ============================================================
% f2014_shaanxi_science_q15b_semicircle.tex
% 2014 陕西卷（理科）第 15 题【B】：以 BC 为直径的半圆
% 来源 PDF：普通高考/2014/2014陕西理.pdf
% 原页：第 1 页（A3），印面 911，第 15 题【B】题干下方
% 图中 BC=6，E、F 在半圆上，AC=2AE；取 b=5、c=(5+√69)/2，
% 使 BC=6 且 ∠A=60°，EF=3。
% ============================================================
\begin{tikzpicture}[>=Stealth, line width=0.7pt, font=\small, x=0.72cm, y=0.72cm]
  \coordinate (B) at (-3,0);
  \coordinate (C) at (3,0);
  \coordinate (O) at (0,0);
  \coordinate (A) at (1.606,4.802);
  \coordinate (E) at (-0.125,2.997);
  \coordinate (F) at (2.533,1.607);

  % 半圆与三角形
  \draw[domain=0:180, samples=120] plot ({3*cos(\x)},{3*sin(\x)});
  \draw (B) -- (A) -- (C) -- cycle;
  \draw (E) -- (F);

  % 点与标签
  \fill (B) circle (1.1pt) node[below left] {$B$};
  \fill (C) circle (1.1pt) node[below right] {$C$};
  \fill (A) circle (1.1pt) node[above] {$A$};
  \fill (E) circle (1.1pt) node[above left] {$E$};
  \fill (F) circle (1.1pt) node[above right] {$F$};
\end{tikzpicture}
\end{center}
\nopagebreak[4]
【C】在极坐标系中，点 $\left(2,\dfrac\pi6\right)$ 到直线 $\rho\sin\left(\theta-\dfrac\pi6\right)=1$ 的距离是\gkblank{}。

\Needspace{4\baselineskip}
\textbf{解析}

\Needspace{5\baselineskip}
\textbf{（A）} 由柯西不等式，
\nopagebreak[4]
\[(ma+nb)^2\leqslant(m^2+n^2)(a^2+b^2).\]
\Needspace{5\baselineskip}
代入 $ma+nb=5$、$a^2+b^2=5$，得
\nopagebreak[4]
\[25\leqslant5(m^2+n^2),\]
\Needspace{5\baselineskip}
故
\nopagebreak[4]
\[m^2+n^2\geqslant5,\qquad \sqrt{m^2+n^2}\geqslant\sqrt5.\]
当 $(m,n)=(a,b)$ 时，$ma+nb=a^2+b^2=5$，等号可以取到，所以最小值为 $\sqrt5$。

\Needspace{5\baselineskip}
\textbf{（B）} 因为 $BC$ 是半圆的直径，所以
\nopagebreak[4]
\[\angle BEC=\angle BFC=90^\circ,\]
即 $CE\perp AB$，$BF\perp AC$，$E$、$F$ 分别是垂足。

\Needspace{5\baselineskip}
在直角三角形 $AEC$ 中，$AE$ 是 $AC$ 在 $AB$ 上的投影，故
\nopagebreak[4]
\[AE=AC\cos A.\]
\Needspace{5\baselineskip}
已知 $AC=2AE$，代入得
\nopagebreak[4]
\[AE=2AE\cos A,\]
\Needspace{5\baselineskip}
而 $AE>0$，所以
\nopagebreak[4]
\[\cos A=\frac12,\qquad A=60^\circ.\]
\Needspace{5\baselineskip}
同理
\nopagebreak[4]
\[AF=AB\cos A=\frac{AB}{2}.\]
\Needspace{5\baselineskip}
设 $AC=b$，$AB=c$，则
\nopagebreak[4]
\[AE=\frac b2,\qquad AF=\frac c2.\]
\Needspace{5\baselineskip}
在 $\triangle AEF$ 中用余弦定理：
\nopagebreak[4]
\[\begin{aligned}
EF^2&=\left(\frac b2\right)^2+\left(\frac c2\right)^2
-2\cdot\frac b2\cdot\frac c2\cos60^\circ\\
&=\frac{b^2+c^2-bc}{4}.
\end{aligned}\]
\Needspace{5\baselineskip}
另一方面，在 $\triangle ABC$ 中由余弦定理，
\nopagebreak[4]
\[BC^2=b^2+c^2-2bc\cos60^\circ=b^2+c^2-bc.\]
\Needspace{5\baselineskip}
已知 $BC=6$，故
\nopagebreak[4]
\[b^2+c^2-bc=36.\]
\Needspace{5\baselineskip}
于是
\nopagebreak[4]
\[EF^2=\frac{36}{4}=9,\qquad EF=3.\]

\Needspace{5\baselineskip}
\textbf{（C）} 极坐标与直角坐标的关系为
\nopagebreak[4]
\[x=\rho\cos\theta,\qquad y=\rho\sin\theta.\]
\Needspace{5\baselineskip}
直线方程化为
\nopagebreak[4]
\[\begin{aligned}
\rho\sin\left(\theta-\frac\pi6\right)
&=\rho\left(\sin\theta\cos\frac\pi6-\cos\theta\sin\frac\pi6\right)\\
&=\frac{\sqrt3}{2}y-\frac12x=1,
\end{aligned}\]
\Needspace{5\baselineskip}
即
\nopagebreak[4]
\[-x+\sqrt3y-2=0.\]
\Needspace{5\baselineskip}
点 $\left(2,\dfrac\pi6\right)$ 的直角坐标为
\nopagebreak[4]
\[x=2\cos\frac\pi6=\sqrt3,\qquad y=2\sin\frac\pi6=1.\]
\Needspace{5\baselineskip}
它到直线的距离为
\nopagebreak[4]
\[d=\frac{|-x+\sqrt3y-2|}{\sqrt{(-1)^2+(\sqrt3)^2}}
=\frac{|-\sqrt3+\sqrt3-2|}{2}=1.\]

\gkanswer{答案：A 为 $\sqrt5$；B 为 $3$；C 为 $1$。A 用柯西不等式；B 由 $AE=AC\cos A$ 得 $\cos A=\dfrac12$，再用余弦定理；C 化极坐标方程为 $-x+\sqrt3y=2$ 后求点到直线距离。}

\end{answer}
\begin{answer}
\gksolhead{16}{}

\textbf{题目}：

$\triangle ABC$ 的内角 $A$，$B$，$C$ 所对的边分别为 $a$，$b$，$c$。

（1） 若 $a$，$b$，$c$ 成等差数列，证明：$\sin A+\sin C=2\sin(A+C)$；

（2） 若 $a$，$b$，$c$ 成等比数列，求 $\cos B$ 的最小值。

\Needspace{4\baselineskip}
\textbf{解析}

\gksubpart{（1）}

\Needspace{5\baselineskip}
因为 $a$，$b$，$c$ 成等差数列，所以
\nopagebreak[4]
\[2b=a+c.\]
\Needspace{5\baselineskip}
由正弦定理
\nopagebreak[4]
\[\frac a{\sin A}=\frac b{\sin B}=\frac c{\sin C}=2R,\]
\Needspace{5\baselineskip}
可设
\nopagebreak[4]
\[a=2R\sin A,\qquad b=2R\sin B,\qquad c=2R\sin C.\]
\Needspace{5\baselineskip}
代入 $2b=a+c$，约去 $2R$，得
\nopagebreak[4]
\[2\sin B=\sin A+\sin C.\]
\Needspace{5\baselineskip}
又因为
\nopagebreak[4]
\[A+B+C=\pi,\]
\Needspace{5\baselineskip}
所以
\nopagebreak[4]
\[B=\pi-(A+C),\qquad \sin B=\sin(A+C).\]
\Needspace{5\baselineskip}
于是
\nopagebreak[4]
\[2\sin(A+C)=\sin A+\sin C,\]
\Needspace{5\baselineskip}
即
\nopagebreak[4]
\[\sin A+\sin C=2\sin(A+C).\]

\gksubpart{（2）}

\Needspace{5\baselineskip}
因为 $a$，$b$，$c$ 成等比数列，所以
\nopagebreak[4]
\[b^2=ac.\]
\Needspace{5\baselineskip}
由余弦定理，
\nopagebreak[4]
\[\begin{aligned}
\cos B&=\frac{a^2+c^2-b^2}{2ac}\\
&=\frac{a^2+c^2-ac}{2ac}\\
&=\frac{a^2+c^2}{2ac}-\frac12.
\end{aligned}\]
\Needspace{5\baselineskip}
由基本不等式
\nopagebreak[4]
\[a^2+c^2\geqslant2ac,\]
\Needspace{5\baselineskip}
得
\nopagebreak[4]
\[\frac{a^2+c^2}{2ac}\geqslant1.\]
\Needspace{5\baselineskip}
因此
\nopagebreak[4]
\[\cos B\geqslant1-\frac12=\frac12.\]
当且仅当 $a=c$ 时取等号，此时 $b^2=ac=a^2$，即 $a=b=c$，三角形为等边三角形。故 $\cos B$ 的最小值为 $\dfrac12$。

\gkanswer{答案（1）：见解析版}

\gkanswer{答案（2）：$\dfrac12$。由 $b^2=ac$ 得 $\cos B=\dfrac{a^2+c^2-ac}{2ac}\geqslant\dfrac12$。}

\end{answer}
\begin{answer}
\gksolhead{17}{}

\textbf{题目}：

四面体 $ABCD$ 及其三视图如图所示，过棱 $AB$ 的中点 $E$ 作平行于 $AD$，$BC$ 的平面分别交四面体的棱 $BD$，$DC$，$CA$ 于点 $F$，$G$，$H$。
\nopagebreak[4]
\begin{center}% ============================================================
% f2014_shaanxi_science_q17_tetrahedron_views.tex
% 2014 陕西卷（理科）第 17 题：四面体 ABCD 及其三视图
% 来源 PDF：普通高考/2014/2014陕西理.pdf
% 原页：第 1 页（A3），印面 911，第 17 题题干下方
% 由三视图得 DA⊥DB、DA⊥DC、DB⊥DC，且 DA=1、DB=DC=2。
% 左图为立体示意图（平行投影，不按比例）；右图为三视图。
% ============================================================
\begin{tikzpicture}[>=Stealth, line width=0.7pt, font=\small,
  line cap=round, line join=round]
  % ---------- 立体示意图 ----------
  \begin{scope}[scale=1.38]
  % 投影：D 为原点，A 在 z 轴，B 在 x 轴，C 在 y 轴。
  \coordinate (D) at (0,0);
  \coordinate (A) at (0,1.15);
  \coordinate (B) at (-1.24,0.96);
  \coordinate (C) at (2.20,0.56);
  \coordinate (E) at (-0.62,1.055); % AB 中点
  \coordinate (F) at (-0.62,0.48);  % BD 中点
  \coordinate (G) at (1.10,0.28);   % DC 中点
  \coordinate (H) at (1.10,0.855);  % CA 中点

  % 可见棱
  \draw (A) -- (B) -- (C) -- cycle;
  \draw (E) -- (H);
  % 不可见棱与截面边
  \draw[dashed] (A) -- (D) -- (B);
  \draw[dashed] (D) -- (C);
  \draw[dashed] (E) -- (F) -- (G) -- (H);

  \fill (A) circle (1.1pt) node[above] {$A$};
  \fill (B) circle (1.1pt) node[below left] {$B$};
  \fill (C) circle (1.1pt) node[right] {$C$};
  \fill (D) circle (1.1pt) node[below] {$D$};
  \fill (E) circle (1.1pt) node[left] {$E$};
  \fill (F) circle (1.1pt) node[below left] {$F$};
  \fill (G) circle (1.1pt) node[below] {$G$};
  \fill (H) circle (1.1pt) node[above right] {$H$};
  \end{scope}

  % ---------- 三视图 ----------
  % 主视图：横向 DB=2，纵向 DA=1
  \begin{scope}[shift={(5.05,0.25)}]
    \draw (0,1) -- (0,0) -- (2,0) -- cycle;
    \draw[<->] (-0.42,0) -- (-0.42,1);
    \node[left] at (-0.42,0.5) {$1$};
    \draw[<->] (0,-0.42) -- (2,-0.42);
    \node[below] at (1,-0.42) {$2$};
    \node[below] at (1,-0.85) {主视图};
  \end{scope}

  % 左视图：横向 DC=2，纵向 DA=1
  \begin{scope}[shift={(7.85,0.25)}]
    \draw (0,1) -- (0,0) -- (2,0) -- cycle;
    \node[below] at (1,-0.45) {左视图};
  \end{scope}

  % 俯视图：横向 DB=2，纵向 DC=2
  \begin{scope}[shift={(5.05,-4.20)}]
    \draw (0,0) -- (0,2) -- (2,2) -- cycle;
    \draw[<->] (-0.42,0) -- (-0.42,2);
    \node[left] at (-0.42,1) {$2$};
    \draw[<->] (0,2.42) -- (2,2.42);
    \node[above] at (1,2.42) {$2$};
    \node[below] at (1,-0.45) {俯视图};
  \end{scope}
\end{tikzpicture}
\end{center}
\nopagebreak[4]
（1） 证明：四边形 $EFGH$ 是矩形；

（2） 求直线 $AB$ 与平面 $EFGH$ 夹角 $\theta$ 的正弦值。

\Needspace{4\baselineskip}
\textbf{解析}

\gksubpart{（1）}

\Needspace{5\baselineskip}
由三视图，主视图是两直角边长分别为 $2$ 和 $1$ 的直角三角形，俯视图是两直角边长均为 $2$ 的直角三角形，左视图同样是两直角边长分别为 $2$ 和 $1$ 的直角三角形。由此可知
\nopagebreak[4]
\[DA\perp DB,\qquad DA\perp DC,\qquad DB\perp DC,\]
\Needspace{5\baselineskip}
且
\nopagebreak[4]
\[DA=1,\qquad DB=DC=2.\]
\Needspace{5\baselineskip}
于是 $DA\perp$ 平面 $BCD$，从而
\nopagebreak[4]
\[AD\perp BC.\]

\Needspace{5\baselineskip}
因为过 $E$ 的平面平行于 $AD$，它与平面 $ABD$ 的交线 $EF$ 满足
\nopagebreak[4]
\[EF\parallel AD.\]
\Needspace{5\baselineskip}
同理，该平面与平面 $ACD$ 的交线 $GH$ 也满足
\nopagebreak[4]
\[GH\parallel AD.\]
\Needspace{5\baselineskip}
所以
\nopagebreak[4]
\[EF\parallel GH.\]

\Needspace{5\baselineskip}
又该平面平行于 $BC$，它与平面 $ABC$ 的交线 $EH$ 满足
\nopagebreak[4]
\[EH\parallel BC,\]
\Needspace{5\baselineskip}
与平面 $BCD$ 的交线 $FG$ 也满足
\nopagebreak[4]
\[FG\parallel BC.\]
\Needspace{5\baselineskip}
所以
\nopagebreak[4]
\[EH\parallel FG.\]

\Needspace{5\baselineskip}
因此四边形 $EFGH$ 是平行四边形。再由 $EF\parallel AD$、$EH\parallel BC$ 及 $AD\perp BC$，得
\nopagebreak[4]
\[EF\perp EH.\]
有一个角为直角的平行四边形是矩形，故四边形 $EFGH$ 是矩形。

\gksubpart{（2）}

\Needspace{5\baselineskip}
以 $D$ 为原点，$\overrightarrow{DB}$、$\overrightarrow{DC}$、$\overrightarrow{DA}$ 的方向分别为 $x$、$y$、$z$ 轴正方向建立空间直角坐标系。由（1）可设
\nopagebreak[4]
\[D(0,0,0),\quad A(0,0,1),\quad B(2,0,0),\quad C(0,2,0).\]
\Needspace{5\baselineskip}
$E$ 是 $AB$ 的中点，故
\nopagebreak[4]
\[E\left(1,0,\frac12\right).\]
\Needspace{5\baselineskip}
平面 $EFGH$ 过 $E$ 且平行于 $AD$（方向 $(0,0,1)$）和 $BC$（方向 $(-2,2,0)$），其法向量可取
\nopagebreak[4]
\[\boldsymbol n=(0,0,1)\times(-2,2,0)=(2,2,0),\]
\Needspace{5\baselineskip}
化简为
\nopagebreak[4]
\[\boldsymbol n=(1,1,0).\]
事实上平面方程为 $x+y=1$。

\Needspace{5\baselineskip}
直线 $AB$ 的方向向量为
\nopagebreak[4]
\[\overrightarrow{AB}=(2,0,-1).\]
\Needspace{5\baselineskip}
设直线 $AB$ 与平面 $EFGH$ 所成角为 $\theta$，则
\nopagebreak[4]
\[\begin{aligned}
\sin\theta
&=\frac{|\overrightarrow{AB}\cdot\boldsymbol n|}
{|\overrightarrow{AB}|\,|\boldsymbol n|}\\
&=\frac{|2\cdot1+0\cdot1+(-1)\cdot0|}
{\sqrt{2^2+0^2+(-1)^2}\sqrt{1^2+1^2}}\\
&=\frac{2}{\sqrt5\cdot\sqrt2}
=\frac{2}{\sqrt{10}}
=\frac{\sqrt{10}}5.
\end{aligned}\]
所以所求正弦值为 $\dfrac{\sqrt{10}}5$。

\gkanswer{答案（1）：见解析版}

\gkanswer{答案（2）：$\dfrac{\sqrt{10}}5$。由三视图得 $DA=1$，$DB=DC=2$ 且两两垂直，平面 $EFGH$ 的法向量为 $(1,1,0)$，与 $\overrightarrow{AB}=(2,0,-1)$ 求线面角。}

\end{answer}
\begin{answer}
\gksolhead{18}{}

\textbf{题目}：

在直角坐标系 $xOy$ 中，已知点 $A(1,1)$，$B(2,3)$，$C(3,2)$，点 $P(x,y)$ 在 $\triangle ABC$ 三边围成的区域（含边界）上。

（1） 若 $\overrightarrow{PA}+\overrightarrow{PB}+\overrightarrow{PC}=\boldsymbol 0$，求 $|\overrightarrow{OP}|$；

（2） 设 $\overrightarrow{OP}=m\overrightarrow{AB}+n\overrightarrow{AC}$（$m,n\in\mathbf R$），用 $x$，$y$ 表示 $m-n$，并求 $m-n$ 的最大值。

\Needspace{4\baselineskip}
\textbf{解析}

\gksubpart{（1）}

\Needspace{5\baselineskip}
由
\nopagebreak[4]
\[\overrightarrow{PA}+\overrightarrow{PB}+\overrightarrow{PC}=\boldsymbol 0,\]
\Needspace{5\baselineskip}
得
\nopagebreak[4]
\[\overrightarrow{OP}=\frac{\overrightarrow{OA}+\overrightarrow{OB}+\overrightarrow{OC}}3.\]
\Needspace{5\baselineskip}
即 $P$ 是 $\triangle ABC$ 的重心。代入三点坐标：
\nopagebreak[4]
\[x_P=\frac{1+2+3}{3}=2,\qquad y_P=\frac{1+3+2}{3}=2.\]
\Needspace{5\baselineskip}
故 $P(2,2)$，
\nopagebreak[4]
\[|\overrightarrow{OP}|=\sqrt{2^2+2^2}=2\sqrt2.\]

\gksubpart{（2）}

\Needspace{5\baselineskip}
由
\nopagebreak[4]
\[\overrightarrow{AB}=(2-1,3-1)=(1,2),\]
\nopagebreak[4]
\[\overrightarrow{AC}=(3-1,2-1)=(2,1),\]
\Needspace{5\baselineskip}
及
\nopagebreak[4]
\[\overrightarrow{OP}=m\overrightarrow{AB}+n\overrightarrow{AC},\]
\Needspace{5\baselineskip}
得
\nopagebreak[4]
\[\begin{cases}
x=m+2n,\\
y=2m+n.
\end{cases}\]
\Needspace{5\baselineskip}
解这个方程组：
\nopagebreak[4]
\[m=\frac{2y-x}{3},\qquad n=\frac{2x-y}{3}.\]
\Needspace{5\baselineskip}
于是
\nopagebreak[4]
\[\begin{aligned}
m-n&=\frac{2y-x}{3}-\frac{2x-y}{3}\\
&=\frac{3y-3x}{3}=y-x.
\end{aligned}\]

\Needspace{5\baselineskip}
$\triangle ABC$ 的三个顶点为 $A(1,1)$、$B(2,3)$、$C(3,2)$。线性函数 $y-x$ 在三角形区域上的最大值必在顶点处取得，计算：
\nopagebreak[4]
\[A:\ 1-1=0;\qquad B:\ 3-2=1;\qquad C:\ 2-3=-1.\]
\Needspace{5\baselineskip}
故
\nopagebreak[4]
\[(m-n)_{\max}=1.\]

\gkanswer{答案（1）：$2\sqrt2$。$P$ 为 $\triangle ABC$ 的重心 $(2,2)$，故 $|\overrightarrow{OP}|=2\sqrt2$。}

\gkanswer{答案（2）：$m-n=y-x$，最大值为 $1$。由 $x=m+2n$，$y=2m+n$ 解出 $m-n=y-x$，在三角形三个顶点处比较得最大值 $1$。}

\end{answer}
\begin{answer}
\gksolhead{19}{}

\textbf{题目}：

在一块耕地上种植一种作物，每季种植成本为 $1000$ 元，此作物的市场价格和这块地上的产量均具有随机性，且互不影响，其具体情况如下表：

\begin{center}
\begin{tabular}{lll}
\toprule
作物产量（kg） & $300$ & $500$ \\
\midrule
概率 & $0.5$ & $0.5$ \\
\bottomrule
\end{tabular}
\end{center}

\begin{center}
\begin{tabular}{lll}
\toprule
作物市场价格（元/kg） & $6$ & $10$ \\
\midrule
概率 & $0.4$ & $0.6$ \\
\bottomrule
\end{tabular}
\end{center}

（1） 设 $X$ 表示在这块地上种植 $1$ 季此作物的利润，求 $X$ 的分布列；

（2） 若在这块地上连续 $3$ 季种植此作物，求这 $3$ 季中至少有 $2$ 季的利润不少于 $2000$ 元的概率。

\Needspace{4\baselineskip}
\textbf{解析}

\gksubpart{（1）}

\Needspace{5\baselineskip}
设作物产量为 $Y$（kg），市场价格为 $P$（元/kg），则一季利润为
\nopagebreak[4]
\[X=YP-1000.\]
$Y$ 与 $P$ 相互独立，按四种组合计算：

\begin{itemize}
\item $Y=300$，$P=6$：
\end{itemize}
\nopagebreak[4]
\[X=300\times6-1000=800,\]
概率 $0.5\times0.4=0.2$；

\begin{itemize}
\item $Y=300$，$P=10$：
\end{itemize}
\nopagebreak[4]
\[X=300\times10-1000=2000,\]
概率 $0.5\times0.6=0.3$；

\begin{itemize}
\item $Y=500$，$P=6$：
\end{itemize}
\nopagebreak[4]
\[X=500\times6-1000=2000,\]
概率 $0.5\times0.4=0.2$；

\begin{itemize}
\item $Y=500$，$P=10$：
\end{itemize}
\nopagebreak[4]
\[X=500\times10-1000=4000,\]
概率 $0.5\times0.6=0.3$。

\Needspace{5\baselineskip}
合并 $X=2000$ 的两种情形，得分布列
\nopagebreak[4]
\[P(X=800)=0.2,\qquad P(X=2000)=0.5,\qquad P(X=4000)=0.3.\]

\gksubpart{（2）}

\Needspace{5\baselineskip}
由（1），单季利润不少于 $2000$ 元的概率为
\nopagebreak[4]
\[p=P(X\geqslant2000)=P(X=2000)+P(X=4000)=0.5+0.3=0.8.\]
\Needspace{5\baselineskip}
设 $3$ 季中利润不少于 $2000$ 元的季数为 $\xi$，则
\nopagebreak[4]
\[\xi\sim B(3,0.8).\]
\Needspace{5\baselineskip}
所求为
\nopagebreak[4]
\[\begin{aligned}
P(\xi\geqslant2)
&=P(\xi=2)+P(\xi=3)\\
&=\mathrm C_3^2(0.8)^2(0.2)+(0.8)^3\\
&=3\times0.64\times0.2+0.512\\
&=0.384+0.512=0.896.
\end{aligned}\]

\gkanswer{答案（1）：$P(X=800)=0.2$，$P(X=2000)=0.5$，$P(X=4000)=0.3$。利润 $X=$ 产量 $\times$ 价格 $-1000$，按独立事件相乘后合并。}

\gkanswer{答案（2）：$0.896$。单季利润不少于 $2000$ 元的概率为 $0.8$，三季中至少两季的概率为 $\mathrm C_3^2(0.8)^2(0.2)+(0.8)^3$。}

\end{answer}
\begin{answer}
\gksolhead{20}{}

\textbf{题目}：

如图，曲线 $C$ 由上半椭圆 $C_1:\dfrac{y^2}{a^2}+\dfrac{x^2}{b^2}=1$（$a>b>0$，$y\geqslant 0$）和部分抛物线 $C_2:y=-x^2+1$（$y\leqslant 0$）连接而成，$C_1$，$C_2$ 的公共点为 $A$，$B$，其中 $C_1$ 的离心率为 $\dfrac{\sqrt3}{2}$。
\nopagebreak[4]
\begin{center}% ============================================================
% f2014_shaanxi_science_q20_ellipse_parabola.tex
% 2014 陕西卷（理科）第 20 题：上半椭圆与抛物线拼接曲线
% 来源 PDF：普通高考/2014/2014陕西理.pdf
% 原页：第 2 页（A3），印面 912，第 20 题题干右侧
% C1: y^2/4+x^2=1 (y≥0)；C2: y=1-x^2 (y≤0)；
% A(-1,0), B(1,0)，P(7/25,48/25)，Q(5/3,-16/9)，直线 PQ 过 B。
% ============================================================
\begin{tikzpicture}[>=Stealth, line width=0.7pt, font=\small, x=1.35cm, y=1.35cm]
  \draw[->] (-2.45,0) -- (2.55,0) node[right] {$x$};
  \draw[->] (0,-2.2) -- (0,2.75) node[above] {$y$};
  \node[below left] at (0,0) {$O$};

  % 上半椭圆 C1
  \draw[domain=-1:1, samples=120]
    plot (\x,{2*sqrt(1-(\x)*(\x))});
  % 抛物线 C2 的 y≤0 两支
  \draw[domain=1:1.75, samples=60] plot (\x,{1-(\x)*(\x)});
  \draw[domain=-1.75:-1, samples=60] plot (\x,{1-(\x)*(\x)});

  % 过 B 的直线
  \coordinate (B) at (1,0);
  \coordinate (P) at (0.28,1.92);
  \coordinate (Q) at (1.6667,-1.7778);
  \draw (P) -- (Q);
  \draw ($(P)!-0.22!(Q)$) -- (P);
  \draw (Q) -- ($(P)!1.10!(Q)$);

  \fill (B) circle (1.1pt) node[above right] {$B$};
  \fill (-1,0) circle (1.1pt) node[above left] {$A$};
  \fill (P) circle (1.1pt) node[right] {$P$};
  \fill (Q) circle (1.1pt) node[right] {$Q$};
\end{tikzpicture}
\end{center}
\nopagebreak[4]
（1） 求 $a$，$b$ 的值；

（2） 过点 $B$ 的直线 $l$ 与 $C_1$，$C_2$ 分别交于点 $P$，$Q$（均异于点 $A$，$B$），若 $AP\perp AQ$，求直线 $l$ 的方程。

\Needspace{4\baselineskip}
\textbf{解析}

\gksubpart{（1）}

$C_1$ 与 $C_2$ 的公共点满足 $y\geqslant0$ 且 $y\leqslant0$，故 $y=0$。

\Needspace{5\baselineskip}
在 $C_1$ 中令 $y=0$：
\nopagebreak[4]
\[\frac{x^2}{b^2}=1,\qquad x=\pm b.\]
\Needspace{5\baselineskip}
在 $C_2$ 中令 $y=0$：
\nopagebreak[4]
\[0=-x^2+1,\qquad x=\pm1.\]
\Needspace{5\baselineskip}
故公共点为 $(\pm b,0)$，与 $(\pm1,0)$ 重合，所以
\nopagebreak[4]
\[b=1.\]
\Needspace{5\baselineskip}
椭圆 $C_1$ 的长半轴为 $a$，短半轴为 $b$，半焦距为 $c$，则
\nopagebreak[4]
\[c^2=a^2-b^2.\]
\Needspace{5\baselineskip}
离心率
\nopagebreak[4]
\[e=\frac ca=\frac{\sqrt3}{2},\]
\Needspace{5\baselineskip}
故
\nopagebreak[4]
\[\frac{c^2}{a^2}=\frac34,\qquad \frac{a^2-b^2}{a^2}=\frac34,\qquad b^2=\frac14a^2.\]
\Needspace{5\baselineskip}
由 $b=1$ 得
\nopagebreak[4]
\[a^2=4,\qquad a=2.\]
\Needspace{5\baselineskip}
所以
\nopagebreak[4]
\[a=2,\qquad b=1,\]
且 $A(-1,0)$，$B(1,0)$。

\gksubpart{（2）}

\Needspace{5\baselineskip}
过 $B(1,0)$ 的直线若斜率不存在，则为 $x=1$，它与 $C_1$ 的另一交点及 $C_2$ 的另一交点都在 $x=1$ 上，此时 $P$ 或 $Q$ 重合于 $B$，不合题意。故设
\nopagebreak[4]
\[l:\ y=k(x-1).\]

\Needspace{5\baselineskip}
先求 $l$ 与 $C_1$ 的非 $B$ 交点 $P$。把直线方程代入
\nopagebreak[4]
\[\frac{y^2}{4}+x^2=1,\]
\Needspace{5\baselineskip}
得
\nopagebreak[4]
\[\frac{k^2(x-1)^2}{4}+x^2=1.\]
\Needspace{5\baselineskip}
令 $u=x-1$，则 $x=1+u$，方程化为
\nopagebreak[4]
\[\frac{k^2u^2}{4}+(1+u)^2=1,\]
\Needspace{5\baselineskip}
即
\nopagebreak[4]
\[\frac{k^2u^2}{4}+2u+u^2=0,\]
\nopagebreak[4]
\[u\left[\left(1+\frac{k^2}{4}\right)u+2\right]=0.\]
\Needspace{5\baselineskip}
$u=0$ 对应点 $B$，另一解为
\nopagebreak[4]
\[u=-\frac{8}{k^2+4}.\]
\Needspace{5\baselineskip}
故
\nopagebreak[4]
\[x_P=1-\frac{8}{k^2+4}=\frac{k^2-4}{k^2+4},\qquad
y_P=k u=-\frac{8k}{k^2+4}.\]

\Needspace{5\baselineskip}
再求 $l$ 与 $C_2$ 的非 $B$ 交点 $Q$。联立
\nopagebreak[4]
\[k(x-1)=1-x^2=-(x-1)(x+1).\]
\Needspace{5\baselineskip}
因 $Q\neq B$，$x\neq1$，约去 $x-1$：
\nopagebreak[4]
\[k=-(x+1),\]
\Needspace{5\baselineskip}
故
\nopagebreak[4]
\[x_Q=-k-1,\qquad y_Q=k(x_Q-1)=k(-k-2)=-k(k+2).\]

\Needspace{5\baselineskip}
计算向量：
\nopagebreak[4]
\[\overrightarrow{AP}=\left(x_P+1,y_P\right)
=\left(\frac{2k^2}{k^2+4},-\frac{8k}{k^2+4}\right),\]
\nopagebreak[4]
\[\overrightarrow{AQ}=\left(x_Q+1,y_Q\right)
=(-k,-k(k+2)).\]
\Needspace{5\baselineskip}
由 $AP\perp AQ$，
\nopagebreak[4]
\[\overrightarrow{AP}\cdot\overrightarrow{AQ}=0.\]
\Needspace{5\baselineskip}
即
\nopagebreak[4]
\[\frac{2k^2}{k^2+4}(-k)+\frac{-8k}{k^2+4}\cdot[-k(k+2)]=0.\]
\Needspace{5\baselineskip}
化简：
\nopagebreak[4]
\[\frac{-2k^3+8k^2(k+2)}{k^2+4}=0,\]
\nopagebreak[4]
\[\frac{6k^3+16k^2}{k^2+4}=0,\]
\nopagebreak[4]
\[2k^2(3k+8)=0.\]
\Needspace{5\baselineskip}
$k=0$ 时 $l$ 为 $x$ 轴，$Q$ 与 $B$ 重合，不合题意，故
\nopagebreak[4]
\[k=-\frac83.\]
\Needspace{5\baselineskip}
于是
\nopagebreak[4]
\[y=-\frac83(x-1),\]
\Needspace{5\baselineskip}
即
\nopagebreak[4]
\[8x+3y-8=0.\]
此时 $P\left(\dfrac7{25},\dfrac{48}{25}\right)$、$Q\left(\dfrac53,-\dfrac{16}9\right)$，均异于 $A$、$B$，且 $AP\perp AQ$ 成立。

\gkanswer{答案（1）：$a=2$，$b=1$。公共点在 $y=0$ 上，由 $C_2$ 得 $x=\pm1$，故 $b=1$；再由 $e=\dfrac{\sqrt3}{2}$ 得 $a=2$。}

\gkanswer{答案（2）：$8x+3y-8=0$。设 $l:y=k(x-1)$，求出 $P$，$Q$ 坐标，由 $\overrightarrow{AP}\cdot\overrightarrow{AQ}=0$ 解得 $k=-\dfrac83$。}

\end{answer}
\begin{answer}
\gksolhead{21}{}

\textbf{题目}：

设函数 $f(x)=\ln(1+x)$，$g(x)=xf'(x)$，$x\geqslant 0$，其中 $f'(x)$ 是 $f(x)$ 的导函数。

（1） 令 $g_1(x)=g(x)$，$g_{n+1}(x)=g(g_n(x))$，$n\in\mathbf N_+$，求 $g_n(x)$ 的表达式；

（2） 若 $f(x)\geqslant ag(x)$ 恒成立，求实数 $a$ 的取值范围；

（3） 设 $n\in\mathbf N_+$，比较 $g(1)+g(2)+\cdots+g(n)$ 与 $n-f(n)$ 的大小，并加以证明。

\Needspace{4\baselineskip}
\textbf{解析}

\gksubpart{（1）}

\Needspace{5\baselineskip}
先求 $g(x)$。由
\nopagebreak[4]
\[f(x)=\ln(1+x),\]
\Needspace{5\baselineskip}
得
\nopagebreak[4]
\[f'(x)=\frac1{1+x}.\]
\Needspace{5\baselineskip}
故
\nopagebreak[4]
\[g(x)=xf'(x)=\frac{x}{1+x}.\]

\Needspace{5\baselineskip}
计算前几项：
\nopagebreak[4]
\[g_1(x)=\frac{x}{1+x},\]
\nopagebreak[4]
\[g_2(x)=g(g_1(x))
=\frac{\dfrac{x}{1+x}}{1+\dfrac{x}{1+x}}
=\frac{x}{1+2x},\]
\nopagebreak[4]
\[g_3(x)=g(g_2(x))
=\frac{\dfrac{x}{1+2x}}{1+\dfrac{x}{1+2x}}
=\frac{x}{1+3x}.\]
\Needspace{5\baselineskip}
猜想
\nopagebreak[4]
\[g_n(x)=\frac{x}{1+nx}.\]
\Needspace{5\baselineskip}
用数学归纳法证明。$n=1$ 时显然成立。假设
\nopagebreak[4]
\[g_n(x)=\frac{x}{1+nx},\]
\Needspace{5\baselineskip}
则
\nopagebreak[4]
\[\begin{aligned}
g_{n+1}(x)
&=g(g_n(x))
=\frac{\dfrac{x}{1+nx}}{1+\dfrac{x}{1+nx}}\\
&=\frac{x}{1+nx+x}
=\frac{x}{1+(n+1)x}.
\end{aligned}\]
\Needspace{5\baselineskip}
故对一切 $n\in\mathbf N_+$，
\nopagebreak[4]
\[g_n(x)=\frac{x}{1+nx}.\]

\gksubpart{（2）}

\Needspace{5\baselineskip}
$f(x)\geqslant ag(x)$ 即
\nopagebreak[4]
\[\ln(1+x)\geqslant a\cdot\frac{x}{1+x},\qquad x\geqslant0.\]
$x=0$ 时两边均为 $0$，对任意 $a$ 成立。

\Needspace{5\baselineskip}
当 $x>0$ 时，$\dfrac{x}{1+x}>0$，不等式等价于
\nopagebreak[4]
\[a\leqslant\frac{(1+x)\ln(1+x)}{x}.\]
\Needspace{5\baselineskip}
记
\nopagebreak[4]
\[h(x)=\frac{(1+x)\ln(1+x)}{x},\qquad x>0.\]
\Needspace{5\baselineskip}
求导：
\nopagebreak[4]
\[\begin{aligned}
h'(x)&=\frac{\left[\ln(1+x)+1\right]x-(1+x)\ln(1+x)}{x^2}\\
&=\frac{x-\ln(1+x)}{x^2}.
\end{aligned}\]
\Needspace{5\baselineskip}
令
\nopagebreak[4]
\[\varphi(x)=x-\ln(1+x),\]
\Needspace{5\baselineskip}
则
\nopagebreak[4]
\[\varphi'(x)=1-\frac1{1+x}=\frac{x}{1+x}>0\qquad(x>0),\]
\Needspace{5\baselineskip}
且 $\varphi(0)=0$，故 $x>0$ 时
\nopagebreak[4]
\[\varphi(x)>0,\]
\Needspace{5\baselineskip}
从而
\nopagebreak[4]
\[h'(x)>0.\]
所以 $h(x)$ 在 $(0,+\infty)$ 上单调递增。

\Needspace{5\baselineskip}
又由洛必达法则或等价无穷小，
\nopagebreak[4]
\[\lim_{x\to0^+}h(x)
=\lim_{x\to0^+}\frac{(1+x)\ln(1+x)}x=1.\]
\Needspace{5\baselineskip}
因此 $h(x)>1$ 对一切 $x>0$ 成立，且下确界为 $1$。要使 $a\leqslant h(x)$ 恒成立，必须且只需
\nopagebreak[4]
\[a\leqslant1.\]
\Needspace{5\baselineskip}
故实数 $a$ 的取值范围为
\nopagebreak[4]
\[(-\infty,1].\]

\gksubpart{（3）}

\Needspace{5\baselineskip}
由（1），
\nopagebreak[4]
\[g(k)=\frac{k}{1+k}=1-\frac1{k+1}.\]
\Needspace{5\baselineskip}
所以
\nopagebreak[4]
\[\begin{aligned}
g(1)+g(2)+\cdots+g(n)
&=\sum_{k=1}^n\left(1-\frac1{k+1}\right)\\
&=n-\sum_{k=1}^n\frac1{k+1}.
\end{aligned}\]
\Needspace{5\baselineskip}
而
\nopagebreak[4]
\[n-f(n)=n-\ln(1+n).\]
\Needspace{5\baselineskip}
因此只需比较
\nopagebreak[4]
\[\sum_{k=1}^n\frac1{k+1}\quad\text{\text{与}}\quad\ln(1+n).\]

\Needspace{5\baselineskip}
对每个 $k=1,2,\cdots,n$，当 $x\in(k,k+1)$ 时，
\nopagebreak[4]
\[\frac1x>\frac1{k+1},\]
\Needspace{5\baselineskip}
两边在 $[k,k+1]$ 上积分得
\nopagebreak[4]
\[\int_k^{k+1}\frac{\mathrm dx}{x}>\int_k^{k+1}\frac{\mathrm dx}{k+1}=\frac1{k+1}.\]
\Needspace{5\baselineskip}
即
\nopagebreak[4]
\[\ln\frac{k+1}{k}>\frac1{k+1}.\]
\Needspace{5\baselineskip}
对 $k=1,2,\cdots,n$ 累加：
\nopagebreak[4]
\[\begin{aligned}
\ln(1+n)
&=\sum_{k=1}^n\ln\frac{k+1}{k}\\
&>\sum_{k=1}^n\frac1{k+1}.
\end{aligned}\]
\Needspace{5\baselineskip}
于是
\nopagebreak[4]
\[\begin{aligned}
g(1)+g(2)+\cdots+g(n)
&=n-\sum_{k=1}^n\frac1{k+1}\\
&>n-\ln(1+n)\\
&=n-f(n).
\end{aligned}\]
\Needspace{5\baselineskip}
所以
\nopagebreak[4]
\[g(1)+g(2)+\cdots+g(n)>n-f(n).\]

\gkanswer{答案（1）：$g_n(x)=\dfrac{x}{1+nx}$。由 $g(x)=\dfrac{x}{1+x}$ 迭代并用归纳法即得。}

\gkanswer{答案（2）：$(-\infty,1]$。$x>0$ 时 $a\leqslant\dfrac{(1+x)\ln(1+x)}x$，该式在 $(0,+\infty)$ 上递增且趋于 $1$。}

\gkanswer{答案（3）：见解析版}

\end{answer}